% Рисунок с доски: область интегрирования D в характеристических переменных.
% Прямая eta = xi отвечает t = 0; треугольник D имеет вершины
% (xi_0, xi_0), (xi_0, eta_0), (eta_0, eta_0).
\begin{figure}[htbp]
\centering
\begin{tikzpicture}[>={Latex[length=4pt]}, scale=1.0]
  \draw[->] (-0.3,0) -- (4.6,0) node[anchor=north east] {$\xi$};
  \draw[->] (0,-0.3) -- (0,4.6) node[anchor=south east] {$\eta$};

  % биссектриса eta = xi : образ прямой t = 0
  \draw[thick, gray] (0,0) -- (4.2,4.2);
  \node[gray, anchor=west, font=\footnotesize] at (3.5,3.1)
    {$t = 0 \iff \xi = \eta$};

  % треугольник D
  \fill[blue!10] (1.2,1.2) -- (1.2,3.4) -- (3.4,3.4) -- cycle;
  \draw[thick, blue!60!black] (1.2,1.2) -- (1.2,3.4) -- (3.4,3.4) -- cycle;
  \node[blue!60!black, font=\footnotesize] at (1.75,2.75) {$D$};

  % угол 45 градусов при нижней вершине
  \draw[gray] (1.2,1.65) arc (90:45:0.45);
  \node[gray, font=\scriptsize] at (1.52,1.78) {$45^\circ$};

  % вершина (xi_0, eta_0)
  \fill (1.2,3.4) circle (1.6pt);
  \node[anchor=south east, font=\footnotesize] at (1.2,3.45) {$(\xi_0,\eta_0)$};

  \draw[dashed, gray] (1.2,0) -- (1.2,1.2);
  \draw[dashed, gray] (0,3.4) -- (1.2,3.4);
  \node[anchor=north, font=\footnotesize] at (1.2,-0.05) {$\xi_0$};
  \node[anchor=east,  font=\footnotesize] at (-0.05,3.4) {$\eta_0$};
  \node[anchor=north, font=\footnotesize] at (3.4,-0.05) {$\eta_0$};
  \draw[dashed, gray] (3.4,0) -- (3.4,3.4);
\end{tikzpicture}
\caption{Область интегрирования $D$ в характеристических переменных:
$\xi_0 \le \xi \le \eta_0$, $\xi \le \eta \le \eta_0$. Две стороны — отрезки
характеристик $\xi = \xi_0$ и $\eta = \eta_0$, третья лежит на прямой
$\xi = \eta$, то есть на оси $t = 0$.}
\label{fig:char-xieta}
\end{figure}
